% from ega4-10-fr.tex
\oldpage[IV-3]{115}
nous tiendrons-nous dans le cadre fixé par Serre. Le présent numéro, tout comme
le précédent, ne sera d'ailleurs pas utilisé dans la suite de ce Traité, et
nous nous bornerons donc à de brèves indications.
\end{env}

\begin{env}[10.10.2]
\label{IV.10.10.2-fr}
Étant donné un ultra-préschéma fixe $R$, on peut naturellement (comme dans
toute catégorie) définir la notion de $R$-\emph{ultra-préschéma}. Considérons
en particulier un corps algébriquement clos $k$~; $\operatorname{Spm}(k)$ est
alors identique à $\operatorname{Spec}(k)$~; nous dirons qu'un
$\operatorname{Spm}(k)$-ultra-préschéma $X'$ est un \emph{espace
préalgébrique sur $k$}~: c'est donc un espace $k$-annelé en anneaux locaux dont
tout point a un voisinage ouvert isomorphe au spectre maximal d'une $k$-algèbre
de type fini~; il revient au même de dire (par (10.9.6)) que $X'=S(X)$, où $X$
est un préschéma localement de type fini sur $k$. Si $X$, $Y$, $Z$ sont trois
$k$-préschémas localement de type fini, il en est de même de
$X\times_ZY$ (\textbf I, 3.4), donc en vertu de (10.9.6), la notion de produit
existe dans la catégorie des espaces préalgébriques sur $k$ (aussi bien le
produit «~sur $k$~» $X'\times_kY'$ que le «~produit fibré~»
$X'\times_{Z'}Y'$, où $X'$, $Y'$, $Z'$ sont trois espaces préalgébriques sur
$k$). On peut donc définir le morphisme diagonal
$\Delta_{X'}:X'\to X'\times_kX'$ (qui n'est autre d'ailleurs que
$S(\Delta_X)$)~; on dit que $X'$ est un \emph{espace algébrique sur $k$} si
$X$ est un schéma, et il revient au même de dire que l'image de $\Delta_{X'}$
est une partie \emph{fermée} de $X'\times_kX'$.
\end{env}

\begin{env}[10.10.3]
\label{IV.10.10.3-fr}
Les simplifications qui proviennent de l'hypothèse que $k$ est algébriquement
clos tiennent d'abord à ce que, pour un $k$-préschéma $X$, localement de type
fini, il y a correspondance biunivoque entre \emph{points fermés de $X$},
\emph{points de $X$ à valeurs dans $k$} (\textbf I, 3.4.4) et \emph{points de
$X$ rationnels sur $k$} (\textbf I, 3.4.5), en vertu de (\textbf I, 6.4.2).
Cela montre en particulier (en vertu de (\textbf I, 3.4.3.1)) que pour deux
espaces $k$-préalgébriques $X'$, $Y'$, l'ensemble sous-jacent au produit
$X'\times_kY'$ est identique à l'\emph{ensemble produit} $X'\times Y'$ des
ensembles sous-jacents (mais bien entendu la topologie de l'espace sous-jacent
à $X'\times_kY'$ n'est pas la \emph{topologie produit} des topologies des
espaces sous-jacents à $X'$ et $Y'$, elle est en général strictement plus fine
que cette dernière).

D'autre part, les anneaux locaux $\mathcal O_x$ aux points d'un espace
préalgébrique $X'$ sur $k$ sont des $k$-algèbres dont on vient de voir que le
corps résiduel est \emph{isomorphe à $k$}~; si $A$ et $B$ sont deux telles
$k$-algèbres locales, tout $k$-homomorphisme $\varphi:A\to B$ est
nécessairement \emph{local}~: en effet, si un élément $x$ de l'idéal maximal de
$A$ était tel que $\varphi(x)$ soit inversible, il existerait $\lambda\in k$
non nul tel que $\varphi(x-\lambda.1)$ appartienne à l'idéal maximal de $B$, ce
qui est absurde puisque $x-\lambda.1$ est inversible dans $A$. On conclut
aussitôt de là que si $X'=S(X)$, $Y'=S(Y)$ sont deux espaces préalgébriques sur
$k$, tout morphisme $X'\to Y'$ d'espaces $k$-annelés est aussi un morphisme
d'espaces $k$-annelés en anneaux locaux (\textbf I, 1.8.2~; cf.
$\mathrm{Err}_{\mathrm{II}}$). D'ailleurs, avec les notations précédentes, si
$A$ et $B$ sont des $k$-algèbres de type fini, $\varphi$ fait de $B$ une
$A$-algèbre de type fini~; donc tout morphisme $X'\to Y'$ de $k$-espaces
annelés est, en vertu de (10.9.6), de la forme $S(f)$, où $f:X\to Y$ est un
morphisme de $k$-préschémas.

Enfin, pour tout ouvert $U$ de $X'$, toute section
$s\in\Gamma(U,\mathcal O_{X'})$, et tout $x\in U$, $s(x)$
($\mathbf 0_{\mathrm I}$, 5.5.1) s'identifie à un élément de $k$, et l'on a
ainsi associé à $s$ une application $\overline s:x\mapsto s(x)$ de $U$ dans
$k$, autrement dit une section au-dessus de $U$ du faisceau $\mathcal A(X')$
des germes d'applications de $X'$ dans $k$~; comme l'application
$h_U:s\mapsto\overline s$ est évidemment un homomorphisme d'anneaux
$\Gamma(U,\mathcal O_{X'})\to\Gamma(U,\mathcal A(X'))$ et commute aux
restrictions à un ouvert $V\subset U$, les $h_U$ définissent un homomorphisme
de faisceaux d'anneaux $h:\mathcal O_{X'}\to\mathcal A(X')$. Si l'on prend
pour $U$ un ouvert ultra-affine $\operatorname{Spm}(A)$, où $A$ est un anneau
de Jacobson, dire que $\overline s=0$ signifie que pour tout idéal maximal
$\mathfrak m$ de $A$, $s$ appartient à $\mathfrak m$, ou encore que $s$ est
dans le \emph{radical} de $A$~; mais comme $A$ est un anneau de Jacobson, son
radical est égal à son nilradical~; pour que $h_U$ soit injective, il faut et
il suffit par suite que $A$ soit réduit.
\end{env}

\begin{env}[10.10.4]
\label{IV.10.10.4-fr}
On dit que l'espace $k$-préalgébrique $X'=S(X)$ est \emph{réduit} s'il en est
ainsi de $X$~; comme l'ensemble des points $x\in X$ où $X$ est réduit est
ouvert ($\mathbf 0_{\mathrm I}$, 5.2.2), son complémentaire contient au moins
un point fermé s'il est non vide (5.1.11), et il revient donc au même de dire
que $X'$ est réduit ou que chacun de ses anneaux locaux $\mathcal O_x$ (pour
$x\in X'$) est réduit. On vient de voir dans (10.10.3) que pour que
l'homomorphisme $h:\mathcal O_{X'}\to\mathcal A(X')$ soit \emph{injectif}, il
faut et il suffit que $X'$ soit réduit. Dans (FAC), Serre se borne en fait aux
espaces préalgébriques \emph{réduits}, ce qui lui permet de définir
$\mathcal O_{X'}$ comme un \emph{sous-faisceau} de $\mathcal A(X')$. On notera
que si $X'$ et $Y'$ sont des $k$-espaces préalgébriques réduits, il en est de
même de $X'\times_kY'$~: en effet, tout revient à voir que si $A$ et $B$ sont
deux $k$-algèbres de type fini réduites, il en est de même de $A\otimes_kB$~;
mais nous avons vu que les radicaux de $A$ et $B$ sont alors réduits à $0$, et
comme $k$ est algébriquement clos, $A$ et $B$ sont des algèbres «~séparables~»
sur $k$ au sens de Bourbaki (Bourbaki, \emph{Alg.}, chap.~VIII, \S~7,
n$^\circ$~5, prop.~5)~; donc $A\otimes_kB$ est sans radical (\emph{loc.
cit.}, n$^\circ$~6, cor.~3 du th.~3), et puisque c'est un anneau de Jacobson,
il est réduit. Toutefois, si $Z'$ est un troisième espace préalgébrique sur
$k$, le «~produit fibré~» $X'\times_{Z'}Y'$ de deux espaces préalgébriques
réduits sur $Z'$ n'est pas réduit en général, ce qui implique que la catégorie
de ces espaces est insuffisante dans de nombreuses questions (notamment dans
la théorie des groupes algébriques). Mais comme on l'a vu ci-dessus, on peut
garder le langage de Serre sans se limiter, comme ce dernier (qui en outre ne
considère que des espaces préalgébriques quasi-compacts), au cas des espaces
préalgébriques réduits.
\end{env}

